\section*{Chapitre \ref{ch:iv1:timed-test-7} --- Épreuve chronométrée 7}
\markboth{\protect\omTocMark{}{Réponses~: Épreuve chronométrée 7}}{\protect\omTocMark{}{Réponses~: Épreuve chronométrée 7}}

\begin{solution}{iq:iv1:timed-test-7:1}
La première vaut 9 lancers, et la seconde 11.

L'attente se découpe en étapes\index{collecte de coupons}. Au départ, l'une ou l'autre des 2 faces recherchées convient, si bien que l'attente vaut $\frac62 = 3$. Il ne reste ensuite qu'une face, et l'attente vaut 6. La somme vaut 9.

Pour trois faces, les étapes donnent $\frac63 + \frac62 + \frac61 = 2 + 3 + 6 = 11$.

Une récurrence exacte et une simulation concordent. La dernière face occupe ici plus de la moitié de l'attente. Elle représente la plus grande part prise isolément dans tout problème de collecte, bien qu'elle devienne inférieure à la moitié dès qu'on en veut quatre ou davantage.

Une seconde méthode, pour la première valeur, utilise l'inclusion--exclusion sur le maximum de deux attentes~: $6 + 6 - 3 = 9$, où 3 est l'attente du premier des deux.

\iqwrong{Additionner les attentes séparées, $6 + 6 = 12$, ce qui compte deux fois les lancers pendant lesquels les deux faces manquent encore.}
\end{solution}

\begin{solution}{iq:iv1:timed-test-7:2}
La probabilité vaut environ $0.554$.

Développons $(q + p)^n$ et $(q - p)^n$\index{astuce de la parité}. Leur somme annule les puissances impaires de $p$, si bien que
\[
P(\text{pair}) = \frac{1 + (1 - 2p)^n}{2} = \frac{1 + 0.8^{10}}{2} = 0.5537.
\]

Cas limites~: pour $p = \frac12$, la formule donne exactement $\frac12$, et pour $p = 0$ elle donne 1, puisque zéro pile est un nombre pair. Quand $n$ croît, elle tend vers $\frac12$ pour tout $p$ strictement compris entre 0 et 1. Une somme binomiale exacte confirme la valeur.

Le même résultat découle d'une récurrence~: un biais de parité $b_n = P(\text{pair}) - P(\text{impair})$ vérifie $b_{n + 1} = (1 - 2p)b_n$, avec $b_0 = 1$.

\iqwrong{Sommer à la main les termes binomiaux pairs en s'arrêtant après quelques-uns, ce qui est lent et n'atteint pas la valeur~: les termes pour 0 et 2 piles ne donnent que $0.54$.}
\end{solution}

\begin{solution}{iq:iv1:timed-test-7:3}
La variance vaut 7, et la pente vaut $\frac12$.

En développant\index{variance d'une somme},
\[
\operatorname{Var}(X + 2Y) = 1 + 4 + 2 \times 2 \times \frac12 = 7.
\]

La pente est $\operatorname{Cov}(Y, X + Y)/\operatorname{Var}(X + Y)$. La covariance vaut $\frac12 + 1 = \frac32$, et la variance vaut $1 + 1 + 1 = 3$. Donc la pente vaut $\frac12$.

Il existe une seconde voie, par symétrie. Les pentes de $X$ et de $Y$ sur $X + Y$ s'additionnent pour donner 1, puisque la régression de $X + Y$ sur elle-même a pour pente 1. Et $X$ et $Y$ ont la même variance, donc les deux pentes sont égales. On obtient ainsi $\frac12$, quelle que soit la corrélation. Une simulation le confirme.

\iqwrong{Oublier un facteur dans le terme croisé. Il vaut $2\operatorname{Cov}(X, 2Y) = 2 \times 2 \times \frac12 = 2$, et une réponse de 6 signifie qu'un facteur 2 a été oublié.}
\end{solution}

\begin{solution}{iq:iv1:timed-test-7:4}
L'espérance vaut $\frac{35}{36} \approx 0.97$.

Par symétrie, $\mathbb{E}\max(X - Y, 0) = \frac12\mathbb{E}|X - Y|$\index{symétrie}. Les écarts de 1, 2, 3, 4 et 5 correspondent à 10, 8, 6, 4 et 2 des 36 issues. Donc
\[
\mathbb{E}|X - Y| = \frac{10 + 16 + 18 + 16 + 10}{36} = \frac{70}{36}.
\]
La moitié vaut $\frac{35}{36}$. Une énumération des 36 issues le confirme.

C'est la valeur d'une très petite option~: on reçoit l'excédent d'un dé sur l'autre, et à la monnaie elle a une valeur positive, bien que $\mathbb{E}(X - Y) = 0$.

\iqwrong{Répondre 0, parce que $X - Y$ est de moyenne nulle. Le maximum écarte les valeurs négatives, si bien qu'une moyenne nulle ne fournit aucun raccourci de ce genre.}

\iqwhy{C'est la forme la plus simple du principe qui sous-tend toute option~: la convexité transforme la variance en valeur.}
\end{solution}

\begin{solution}{iq:iv1:timed-test-7:5}
La première vaut $\frac{1}{11}$, et la seconde $\frac16$.

Sur les 36 issues, 11 comportent au moins un 6, et une seule en comporte deux\index{probabilité conditionnelle}. Cela donne $\frac{1}{11}$.

Si le premier dé est connu, l'autre est indépendant, et la probabilité vaut $\frac16$.

La différence tient à ce que l'on sait~: «~au moins un~» conserve 11 issues équiprobables, et «~celui-ci~» en conserve 6. Une énumération aboutit au même résultat.

Quelle lecture s'applique dépend de la manière dont l'information est arrivée. Si quelqu'un demande «~y a-t-il un 6~?~», c'est la première qui vaut. Si l'on a vu un dé, c'est la seconde. En pratique, cette étape tranche la réponse, et il vaut la peine de la formuler à voix haute.

\iqwrong{Répondre $\frac16$ dans les deux cas, au motif que l'autre dé est indépendant. Dans le premier cas, il n'y a pas d'«~autre dé~» à désigner.}
\end{solution}

\begin{solution}{iq:iv1:timed-test-7:6}
La probabilité vaut $\frac38$.

Le déterminant $ad - bc$ est nul lorsque $ad = bc$\index{déterminant}. Chaque produit vaut 1 dans un cas sur 4.

Les deux produits valent 0 dans $3 \times 3 = 9$ des 16 matrices, et valent tous deux 1 dans une matrice de plus. Il y a donc 10 matrices singulières, et 6 inversibles, ce qui donne $\frac{6}{16}$.

Une énumération des 16 cas donne le même résultat. Sur les réels, les 6 sont celles où exactement l'un des deux produits vaut 1. Quand la taille croît, la probabilité de singularité pour des matrices binaires aléatoires tend vers 0, si bien que cette faible probabilité n'est qu'un effet de petite taille.

\iqwrong{Supposer que les matrices aléatoires sont presque toujours inversibles. C'est vrai pour des coefficients continus. Avec des coefficients binaires, la coïncidence $ad = bc$ est fréquente.}
\end{solution}

\begin{solution}{iq:iv1:timed-test-7:7}
La probabilité vaut $\frac{99}{116} \approx 0.853$.

Les vraisemblances de deux piles valent $\frac14$, 1 et $\frac{9}{16}$\index{formule de Bayes}. Multipliées par 16, elles valent 4, 16 et 9, soit un total de 29. Les poids a posteriori sont donc $\frac{4}{29}$, $\frac{16}{29}$ et $\frac{9}{29}$.

Le lancer suivant donne pile avec la probabilité
\[
\frac{4 \cdot \frac12 + 16 \cdot 1 + 9 \cdot \frac34}{29} = \frac{24.75}{29} = \frac{99}{116}.
\]
Une simulation confirme. Avant tout lancer, la probabilité valait $\frac{2.25}{3} = \frac34$. Deux piles la portent à $0.853$, un mouvement raisonnable, car l'observation favorise la pièce à deux piles sans pouvoir écarter les deux autres.

\iqwrong{Retenir la pièce la plus vraisemblable et répondre 1. La loi a posteriori est une pondération des trois pièces.}

\iqfollow{Et si les deux lancers donnaient P, F~?}{La pièce à deux piles est exclue. Les poids valent $\frac14$ et $\frac{3}{16}$, soit 4 et 3, et le lancer suivant donne pile avec la probabilité $\frac{4 \cdot \frac12 + 3 \cdot \frac34}{7} = \frac{17}{28}$.}
\end{solution}

\enlargethispage{-4\baselineskip}%
\begin{solution}{iq:iv1:timed-test-7:8}
L'espérance de l'âge vaut $\frac23$, et l'espérance de la durée de vie totale vaut $\frac43$.

Un instant tiré au hasard tombe plus volontiers dans une longue durée de vie. Une ampoule de durée de vie $\ell$ est observée avec un poids proportionnel à $\ell$\index{paradoxe de l'inspection}. Ainsi, l'ampoule en service a une espérance de durée de vie de
\[
\frac{\mathbb{E}L^2}{\mathbb{E}L} = \frac{4/3}{1} = \frac43,
\]
ce qui dépasse la moyenne simple de 1.

Dans la vie de cette ampoule, l'instant est uniforme, donc l'espérance de l'âge vaut la moitié de la durée de vie~: $\frac23$. Et l'espérance de la durée de vie restante vaut elle aussi $\frac23$. Une simulation d'une longue série d'ampoules le confirme.

\iqwrong{Répondre 1 et $\frac12$. Ces réponses portent sur une ampoule typique, mais un échantillonnage à un instant aléatoire privilégie les plus longues.}

\iqwhy{Le même effet fait qu'un passager attend plus longtemps que la moitié de l'intervalle moyen entre deux bus.}
\end{solution}

\begin{solution}{iq:iv1:timed-test-7:9}
$\mathbb{E}[Z \mid Z > 1] \approx 1.52$, et l'expected shortfall vaut environ $2.67$.

Pour une normale centrée réduite, $\varphi^\prime(z) = -z\varphi(z)$, ce qui donne $\int_a^\infty z\varphi(z)\,dz = \varphi(a)$. Donc\index{loi normale tronquée}
\[
\mathbb{E}[Z \mid Z > a] = \frac{\varphi(a)}{P(Z > a)}.
\]
En $a = 1$, cela vaut $\frac{0.242}{0.159} = 1.52$.

L'expected shortfall est la même formule évaluée au quantile~: $\frac{0.0267}{0.01} = 2.67$. Il n'est ici qu'environ $15~\%$ au-dessus de la valeur en risque, car la queue décroît vite. Une queue plus épaisse donne un écart plus grand. Une intégration numérique confirme les deux valeurs.

Un contrôle en $a = 0$~: la formule donne $\frac{0.399}{0.5} = 0.80$, qui est bien la moyenne de $|Z|$, comme attendu.

\iqwrong{Diviser par la mauvaise queue, ou donner le quantile lui-même. La réponse doit être supérieure à $a$, et ce contrôle détecte la première erreur.}
\end{solution}

\begin{solution}{iq:iv1:timed-test-7:10}
La moyenne vaut 3, et la variance vaut $\frac{13}{6}$.

Chacune des $\binom42 = 6$ paires est inversée avec probabilité $\frac12$. La moyenne vaut donc 3\index{variables indicatrices}.

Pour la variance, construire la permutation en insérant les éléments dans l'ordre. L'élément $k$ prend l'une des $k$ positions et ajoute un nombre uniforme sur $\{0, \dots, k - 1\}$ d'inversions, indépendant des précédents. Une loi uniforme sur $\{0, \dots, k - 1\}$ a pour variance $\frac{k^2 - 1}{12}$. Donc
\[
\operatorname{Var} = \frac{0 + 3 + 8 + 15}{12} = \frac{26}{12} = \frac{13}{6}.
\]
Une énumération des 24 permutations le confirme. Le résultat général est $n(n - 1)(2n + 5)/72$. L'étape d'insertion supprime toutes les covariances de la voie des indicatrices.

\iqwrong{Additionner les variances des paires, $6 \times \frac14 = \frac32$. Les paires qui partagent une position ne sont pas indépendantes.}
\end{solution}

\begin{solution}{iq:iv1:timed-test-7:11}
L'espérance vaut $\sqrt{\frac{0.5}{\pi}} \approx 0.399$.

Écrivons $\max(X, Y) = \frac{X + Y}{2} + \frac{|X - Y|}{2}$\index{maximum et valeur absolue}. Le premier terme est de moyenne nulle. La différence $X - Y$ est normale de variance $2 - 2\rho = 1$.

Pour une normale de variance $s^2$, $\mathbb{E}|\cdot| = s\sqrt{2/\pi}$. La moyenne du maximum vaut donc $\frac12\sqrt{2/\pi}\sqrt{2 - 2\rho} = \sqrt{(1 - \rho)/\pi}$. En $\rho = \frac12$, on obtient $0.399$.

Contrôles aux limites~: en $\rho = 1$, les deux variables sont égales et le maximum a une moyenne nulle. En $\rho = 0$, il vaut $\frac{1}{\sqrt\pi} = 0.564$, valeur connue pour deux normales indépendantes. Une simulation donne $0.399$.

\iqwrong{Multiplier la réponse indépendante par $1 - \rho$, ce qui donne $0.282$. La dépendance passe par $\sqrt{1 - \rho}$.}
\end{solution}

\begin{solution}{iq:iv1:timed-test-7:12}
La probabilité vaut environ $0.04$.

La somme a pour moyenne 350 et pour variance $100 \cdot \frac{35}{12} = 291.7$, donc son écart type vaut environ $17.1$\index{théorème central limite}.

Avec la correction de continuité, le seuil vaut $379.5$, et $z = \frac{29.5}{17.1} = 1.73$. Ainsi $P \approx 0.042$. La somme de 100 dés est assez proche d'une normale à ce niveau, et une convolution exacte donne la même valeur à $0.001$ près.

Sans la correction, $z = 1.75$ et la queue vaut $0.040$. Pour une réponse rapide, cela suffit, même si la correction coûte une étape supplémentaire pour le chiffre en plus.

\iqwrong{Prendre une variance de 25, comme pour 100 lancers de pièce. La variance d'un dé est proche de 3, ce qui est un repère utile à connaître.}
\end{solution}
